\subsection{Summarization and Coarsening}
\label{sec:method:reduction:summarization}
A coarsening is a quotient under a merge map $\mu \colon \mathcal{V} \rightarrow \mathcal{V}'$
(\ref{sec:prelim:reduction:summarization}), and, unlike the other two families, which only
delete, must also decide what a super-node \emph{is}. The rewrite that does so is fixed once
(\ref{sec:method:summary:mergemap}), and every candidate below reduces to its choice of
$\mu$.
% Original: A coarsening is a quotient. A merge map $\mu \colon \mathcal{V} \rightarrow \mathcal{V}'$
% sends each node to the super-node containing it, and the reduced graph is the quotient of $G$
% under $\mu$ (\ref{sec:prelim:reduction:summarization}). Unlike the other two families, which
% only delete, a coarsening must also decide what a super-node \emph{is}. The rewrite that
% turns a merge map into a graph is therefore fixed once
% (\ref{sec:method:summary:mergemap}), and every candidate below reduces to its choice of
% $\mu$.

Four candidates span domain-specific, exact and general-purpose, with a random control
making matched-compression comparison possible (Table~\ref{tab:summarization_methods}). No
summarization configuration has been trained, so what follows states the intended
construction, not a measured method.
% Original: Four candidates span domain-specific, exact and general-purpose, with a random control that
% is what makes matched-compression comparison possible
% (Table~\ref{tab:summarization_methods}). No summarization configuration has been trained, so
% what follows states the intended construction rather than a measured method.

\begin{table}[htbp]
\centering
\caption{The summarization candidates, the equivalence each merges by, and whether its
compression can be set. Only a method taking a target ratio can be matched to another
directly; the rest are matched through the random control
of~\ref{sec:method:summary:random}.}
\label{tab:summarization_methods}
\small
\begin{tabular}{llll}
\toprule
\textbf{Method} & \textbf{Merge criterion} & \textbf{Role} & \textbf{Compression} \\
\midrule
Cone or MFFC contraction & fanout-free cones        & domain-specific & graph-dependent \\
Graded colour refinement & $d$-round colour classes & exact anchor    & graph-dependent \\
Convolution matching     & convolution equivalence  & general bar     & target ratio \\
Random within-type       & uniform within type      & control         & exact \\
\bottomrule
\end{tabular}
\end{table}

\subsubsection{Merge Map and Super-Node Schema}
\label{sec:method:summary:mergemap}
One rewrite turns any merge map into a graph, so methods differ only in what they merge. A
super-node carries the counts of its members by role, a super-edge the counts of normal and
inverted connections it stands for,
\begin{equation}
    x_{u'} = \Bigl( \bigl| \{\, u \in \mu^{-1}(u') : \tau(u) = r \,\} \bigr|
        \Bigr)_{r \in \{\mathrm{const},\, \mathrm{PI},\, \mathrm{AND},\, \mathrm{PO}\}},
    \qquad
    a_{u'v'} = \sum_{\substack{(u,v) \in \mathcal{E} \\ \mu(u) = u',\ \mu(v) = v'}}
        \begin{pmatrix} 1 - \mathrm{inv}(u,v) \\ \mathrm{inv}(u,v) \end{pmatrix},
    \label{eq:method:summary-rewrite}
\end{equation}
with $\tau$ the role of~\eqref{eq:prelim:role} and $\mathrm{inv}$ the inversion
of~\eqref{eq:prelim:inversion}. Super-edges are kept only for $u' \neq v'$. Intra-cluster
edges become self-loops and are dropped, their count retained as a graph attribute. The two
entries of $a_{u'v'}$ sum to the multiplicity $w_{u'v'}$ of~\eqref{eq:prelim:multiplicity},
and positional encodings are pooled by member minimum, so a super-node sits at the level of
its earliest member.
% TODO: super-node content is not yet exploited; do not write as though it is.
% Original: One rewrite turns any merge map into a graph, so that methods differ in what they merge and
% in nothing else. A super-node carries the counts of its members by role, and a super-edge the
% counts of the normal and inverted connections it stands for,
% [[equation unchanged, see above]]
% with $\tau$ the role of~\eqref{eq:prelim:role} and $\mathrm{inv}$ the inversion
% of~\eqref{eq:prelim:inversion}. Super-edges are kept only for $u' \neq v'$, so intra-cluster
% edges become self-loops and are dropped; their count is retained as a graph attribute rather
% than lost silently. The two entries of $a_{u'v'}$ sum to the multiplicity $w_{u'v'}$
% of~\eqref{eq:prelim:multiplicity}. Positional encodings are pooled by member minimum, so a
% super-node sits at the level of its earliest member.

\paragraph{Boundary preservation.}
Every candidate merges within one node role,
\begin{equation}
    \mu(u) = \mu(v) \implies \tau(u) = \tau(v).
    \label{eq:method:summary-boundary}
\end{equation}
The PI/PO interface is the circuit's contract with the rest of the design. Dissolving it
changes what the graph represents rather than coarsening it, and imposing the restriction
uniformly makes a comparison between methods a comparison of \emph{which} nodes are merged,
not of whether the interface survived.
% Original: The PI/PO interface is the circuit's contract with the rest of the design, and dissolving it
% changes what the graph represents rather than coarsening it. Imposing the restriction
% uniformly also makes a comparison between methods a comparison of \emph{which} nodes are
% merged, not of whether the interface survived.

\paragraph{Schema compatibility.}
\label{sec:method:summary:schema}
The count representation of~\eqref{eq:method:summary-rewrite} is a strict superset of the
one-hot one: a super-node with a single member reproduces its original feature row exactly,
and a single edge its own inversion indicator. An unreduced AIG is therefore already a valid
input in the summarized schema, without conversion, which is what makes RQ5 testable with
one set of weights. It is a design constraint on every method, not a convenience: a merge
producing features outside this space would leave summary and full graphs in different
input spaces and make cross-state inference undefined.
% Original: The count representation of~\eqref{eq:method:summary-rewrite} is a strict superset of the
% one-hot one: a super-node with a single member reproduces its original feature row exactly,
% and a single edge reproduces its own inversion indicator. An unreduced AIG is therefore
% already a valid input in the summarized schema, without conversion, and that is what makes
% RQ5 testable with one set of weights. It is a design constraint on every method rather than a
% convenience. A merge producing features outside this space would leave summary and full
% graphs in different input spaces and make cross-state inference undefined.

\subsubsection{Domain-Specific Coarsening}
\label{sec:method:summary:domain}
Two AIG-native candidates. The choice between them awaits measured compression on this
corpus.

\emph{MFFC contraction} sends every gate to the root of the maximal fanout-free cone
containing it, $\mu(u) = v$ for the $v$ with $u \in \operatorname{MFFC}(v)$
of~\eqref{eq:prelim:mffc}, iterated to a fixed point. What survives is the multi-fanout
sharing skeleton plus the fixed interface. It takes no parameter, so its compression is a
property of the circuit.
% Original: \emph{MFFC contraction} sends every gate to the root of the maximal fanout-free cone
% containing it, $\mu(u) = v$ for the $v$ with $u \in \operatorname{MFFC}(v)$
% of~\eqref{eq:prelim:mffc}, iterated to a fixed point. What survives is the multi-fanout
% sharing skeleton together with the fixed interface. It takes no parameter, so its compression
% is a property of the circuit.

\emph{Cone coarsening} merges along two axes: chains of single-fanout gates contract into
one super-node up to a capped chain length, reducing depth, and gates inside a bounded band
of topological levels sharing an immediate post-dominator (\ref{sec:prelim:aig:properties})
merge together, compressing width while leaving critical-path depth, and with it the level
encoding, exact. The chain cap and band width are the only compression controls a
domain-specific candidate offers here, and neither has been given a value.
% TODO: fix the chain cap and the level-band width, or drop the cone candidate.
% TODO: report that the width axis must use post-dominators, if this method is written up.
% TODO: neither cone nor MFFC compression is measured on the corpus.
% Original: \emph{Cone coarsening} merges along two axes. Chains of single-fanout gates contract into one
% super-node up to a capped chain length, which reduces depth. Gates inside a bounded band of
% topological levels that share an immediate post-dominator
% (\ref{sec:prelim:aig:properties}) merge together, which compresses width while leaving
% critical-path depth, and with it the level encoding, exact. The chain cap and the band width
% are the only compression controls a domain-specific candidate offers here, and neither has
% been given a value.

The two are related, not rival: chain contraction is MFFC absorption's capped, chain-shaped
special case, which additionally captures reconverging cone interiors. Both preserve
acyclicity, since a cluster absorbs only into the cluster receiving all its outgoing edges,
which cannot close a cycle in a DAG\@. For cone coarsening this holds only at band width
zero, and wider bands lose the level encoding's exactness too.
% Original: The two are related rather than rival. Chain contraction is the capped, chain-shaped special
% case of MFFC absorption, which additionally captures cones whose interior reconverges. Both
% preserve acyclicity, since a cluster is absorbed only into the single cluster receiving all
% of its outgoing edges, which cannot close a cycle in a DAG\@. For cone coarsening the
% guarantee holds at band width zero, where the width axis merges only within a level; wider
% bands give it up along with the exactness of the level encoding.

Merging MFFCs invites an objection, and the answer is the thesis claim in miniature:
refactoring collapses and re-expresses one MFFC at a time, so merging them might be expected
to destroy exactly the signal the model must read. The claim here is the opposite, and
falsifiable: intra-cone wiring is what local rewriting is \emph{free} to change, hence noise
with respect to the label, while the sharing structure that \emph{constrains} rewriting is
what survives the merge.
% Original: Merging MFFCs invites an objection, and the answer is the thesis claim in miniature.
% Refactoring collapses and re-expresses one MFFC at a time, so merging them might be expected
% to destroy exactly the signal the model must read. The claim here is the opposite, and it is
% falsifiable: intra-cone wiring is what local rewriting is \emph{free} to change, so it is
% noise with respect to the label, whereas the sharing structure that \emph{constrains} what
% rewriting can achieve is what survives the merge.

\subsubsection{Graded Colour Refinement}
\label{sec:method:summary:wl}
The exact anchor of the study: the one candidate whose merge map is defined by what the
encoder can see, not by a property of the graph. Nodes merge when they agree after $d$
rounds of colour refinement~\eqref{eq:prelim:cr} (\ref{sec:prelim:gnn:expressivity}), and two
substitutions in that recurrence carry the whole adaptation to AIGs,
\begin{equation}
    c_v^{(0)} = \bigl( \tau(v),\ \ell(v) \bigr),
    \qquad
    \bigl\{\!\!\bigl\{\, c_u^{(t-1)} : u \in \mathcal{N}(v) \,\bigr\}\!\!\bigr\}
    \;\longmapsto\;
    \bigl\{\!\!\bigl\{\, \bigl( c_u^{(t-1)},\ \mathrm{inv}(u,v) \bigr) :
        u \in \operatorname{fi}(v) \,\bigr\}\!\!\bigr\}_{c}.
    \label{eq:method:summary-refine}
\end{equation}
Initialising from role and level~\eqref{eq:prelim:level} rather than uniformly keeps nodes
at different depths in different classes, so a super-node's pooled level is exact. Carrying
$\mathrm{inv}$ inside the neighbour's identity keeps a normal and an inverted connection out
of the same class. Without it, logically distinct subgraphs merge. The neighbourhood is the
fanin~\eqref{eq:prelim:fanin}, so refinement runs backward from the primary outputs, itself
a parameter that changes both how much compresses and which structure survives.
% TODO: fix the refinement direction, or run both and report each.
% Original: The exact anchor of the study, and the one candidate whose merge map is defined by what the
% encoder can see rather than by a property of the graph. Nodes merge when they agree after $d$
% rounds of the refinement~\eqref{eq:prelim:cr}, and two substitutions in that recurrence carry
% the whole adaptation to AIGs,
% [[equation unchanged, see above]]
% Initialising from the role and the level~\eqref{eq:prelim:level} rather than uniformly keeps
% nodes at different depths in different classes, so a super-node's pooled level is exact.
% Carrying $\mathrm{inv}$ inside the neighbour's identity keeps a normal and an inverted
% connection out of the same class; without it, logically distinct subgraphs merge. The
% neighbourhood is the fanin~\eqref{eq:prelim:fanin}, so refinement runs backward from the
% primary outputs. Direction is itself a parameter, and it changes both how much compresses and
% which structure survives.

The subscript in~\eqref{eq:method:summary-refine} is the count cap
of~\ref{sec:prelim:gnn:expressivity}. Only its two named endpoints, $c = \infty$ and
$c = 1$, are run, since intermediate caps are bracketed by them and each value costs a
training run. The depth is set to $d = L$, the encoder's message-passing depth, and that
equality is the justification, not a convenience: a smaller $d$ merges less and gains
nothing, while a larger one merges nodes the encoder can separate, and the losslessness
argument fails.
% Original: The subscript in~\eqref{eq:method:summary-refine} is the count cap
% of~\ref{sec:prelim:gnn:expressivity}, and only its two named endpoints are run:
% $c = \infty$, lossless for a summing GNN, and $c = 1$, which is bisimulation, coarser and
% lossy for a counting model. Intermediate caps are omitted because the endpoints bracket them
% and each value costs a training run. The depth is set to $d = L$, the encoder's
% message-passing depth, and that equality is the justification rather than a convenience. A
% smaller $d$ merges less and gains nothing; a larger one merges nodes the encoder can
% separate, and the losslessness argument fails.

At $c = \infty$ the merge is orthogonal to the optimization being predicted by construction
(\ref{sec:relwork:reduction:summarization}): the positive control for RQ5, realised end to
end only through the encoder variant of~\ref{sec:method:architecture:exact}.
% Original: At $c = \infty$ the merge is orthogonal to the optimization being predicted by construction:
% it removes only what the model provably cannot see, so it cannot erase the label. That is
% what makes it the positive control for RQ5 rather than merely another contender, and
% realising the guarantee end to end requires the encoder variant
% of~\ref{sec:method:architecture:exact}.

Structural hashing has already removed one-hop equivalences
(\ref{sec:prelim:algorithms:primitives}), since every corpus graph is strashed, so what this
method finds is \emph{residual} redundancy beyond one hop. How much exists is unknown, and
the answer bounds what the exact track can compress.
% TODO: run the residual-redundancy probe over depths 1 to 4; depth 1 subsumes k-SNAP \citep{tian2008ksnap}.
% Original: Structural hashing has already removed the one-hop equivalences
% (\ref{sec:prelim:algorithms:primitives}), since every corpus graph is strashed. What this
% method finds is therefore \emph{residual} redundancy beyond one hop. How much of it exists is
% not yet known, and the answer bounds what the exact track can compress.

\subsubsection{Convolution Matching}
\label{sec:method:summary:convmatch}
ConvMatch \citep{dickens2024convmatch} (\ref{sec:relwork:reduction:summarization}) is
adapted only by imposing the boundary constraint~\eqref{eq:method:summary-boundary} on its
candidate set, leaving the convolution objective itself unmodified. It accepts a target
compression ratio directly, and is the bar the domain-specific candidate must clear.
% TODO: set the target ratio from the domain-specific candidate's achieved compression.
% Original: The strongest general-purpose competitor. ConvMatch \citep{dickens2024convmatch} merges nodes
% that are equivalent with respect to the \emph{graph convolution operation} itself rather than
% with respect to a graph property, and it accepts a target compression ratio directly. The
% only adaptation is~\eqref{eq:method:summary-boundary}, imposed on the candidate set so that
% the convolution objective is left unmodified. It is the bar the domain-specific candidate
% must clear, and the right bar because it is GNN-aware and domain-blind: a win against it is a
% win against a method optimising the same objective, not against a strawman.

Its behaviour on a deep DAG with inverted edges is untested in the literature, and one
property of its objective suggests why it may differ here: the merge cost is an unweighted
sum over the convolution output, so merging two low-degree nodes scores better than merging
two genuine convolution twins of higher degree, regardless of similarity. Convolution
equivalence therefore decides \emph{between} pairs of comparable degree, not across them,
which on an AIG's degree profile biases merging toward the input frontier.
% TODO: measure the per-graph cost of probe sampling before quoting it.
% Original: Its behaviour on a deep DAG with inverted edges is untested in the literature, and one
% property of its objective explains why it may differ from the published setting. The merge
% cost is an unweighted sum over the convolution output, so merging two low-degree nodes scores
% better than merging two genuine convolution twins of higher degree, regardless of similarity.
% Convolution equivalence therefore decides \emph{between} pairs of comparable degree and not
% across them, which on the degree profile of an AIG biases merging toward the input frontier.

\subsubsection{Random Within-Type Merging}
\label{sec:method:summary:random}
Uniformly random nodes of the same type are merged until the node count reaches a target,
itself not free: it is set per pairing to the partner method's achieved compression, under
the same boundary restriction~\eqref{eq:method:summary-boundary} as every other candidate
(\ref{sec:method:summary:mergemap}).
% Original: Uniformly random nodes of the same type are merged until the node count reaches a target. The
% target is not a free parameter: it is set per pairing to the partner method's achieved
% compression. Restricting merges by~\eqref{eq:method:summary-boundary} is what makes the
% comparison isolate \emph{which} nodes are merged rather than confounding it with whether the
% interface survived.

As a naive floor, it answers the question a reader asks first: a method that does not beat
random merging at the same compression bought nothing with its sophistication, the
counterpart of random edge dropout (\ref{sec:method:sparse:dropout}) so both halves of the
study share a floor built the same way.
% Original: As a naive floor, it answers the question a reader asks first: a method that does not beat
% random merging at the same compression bought nothing with its sophistication. It is also the
% counterpart of random edge dropout (\ref{sec:method:sparse:dropout}), so both halves of the
% study have a floor built the same way.

As a matching instrument, it is what makes RQ4 answerable for this family
(\ref{sec:relwork:reduction:summarization}): convolution matching takes a target ratio,
cone coarsening offers coarse integer knobs, and MFFC contraction and colour refinement have
none whatsoever, so each is compared against a random control at its own achieved
compression.
% TODO: not built. The cost is the training runs, one per matched point, not the code.
% Original: As a matching instrument, it is what makes RQ4 answerable for this family at all. Compression
% is a free parameter for only some methods: convolution matching takes a target ratio, cone
% coarsening offers coarse integer knobs, and MFFC contraction and colour refinement have none
% whatsoever. There is no general way to make two \emph{real} methods meet at one ratio. Each
% method is instead compared against a random control run at that method's own achieved
% compression, which replaces a matching problem that has no solution.

% Superseded draft, retained for reference until the summarization runs exist.
% \subsection{Summarization and Coarsening}
% \label{sec:method:reduction:summarization}
% These methods compress the graph by merging nodes into super-nodes. Unlike the other two
% families, which only delete, summarization must decide what a merged node \emph{is}, so
% the shared rewrite that turns a clustering into a graph is described first
% (\ref{sec:method:summary:mergemap}), and each method below then reduces to the single
% question of which nodes it groups together.
%
% Three methods are run, spanning domain-specific, general state of the art, and exact, so
% that RQ4 has something to compare. A fourth, random merging, is a control rather than a
% contender; it is what makes matched-compression comparison possible.
% Table~\ref{tab:summarization_methods} lists the set.
%
% \begin{table}[htbp]
% \centering
% \caption{Summarization methods.}
% \label{tab:summarization_methods}
% \small
% \begin{tabular}{lll}
% \toprule
% \textbf{Method} & \textbf{Role} & \textbf{Compression} \\
% \midrule
% Domain-specific coarsening (TBD) & domain-specific        & method-dependent \\
% Colour refinement                & adapted, exact         & fixed \\
% ConvMatch                        & GNN-aware SOTA         & target ratio \\
% Random within-type               & control / matching     & exact \\
% \bottomrule
% \end{tabular}
% \end{table}
%
% \subsubsection{Merge Map and Super-Node Schema}
% \label{sec:method:summary:mergemap}
% Every method produces only a \emph{clustering}: an assignment of each node to a group. One
% shared rewrite turns that clustering into a graph, so that methods differ in what they merge
% and in nothing else. Given a clustering, the rewrite:
%
% \begin{itemize}
%     \item sets each super-node's features to the \emph{member type counts}
%     $[\#\mathrm{const}, \#\mathrm{PI}, \#\mathrm{AND}, \#\mathrm{PO}]$;
%     \item remaps every edge to its endpoints' clusters, discards the resulting self-loops,
%     and sums the inversion counts of parallel super-edges so that a super-edge records how
%     many normal and how many inverted connections it stands for;
%     \item pools the positional encoding by member \emph{minimum}, so a super-node sits at
%     the level of its earliest member;
%     \item records the number of discarded intra-cluster edges, and the preserved primary
%     input and output counts, as graph attributes.
% \end{itemize}
%
% \paragraph{Schema compatibility.}
% \label{sec:method:summary:schema}
% The count representation is a strict superset of the one-hot representation: a super-node
% with one member reproduces its original feature row exactly. An unreduced AIG is therefore
% already a valid input in the summarized schema, without conversion. This is the property
% that makes RQ5 testable with a single set of weights, and it is a design constraint on every
% method rather than a convenience: a merge that produced features outside this space would
% put summary and full graphs in different input spaces and make cross-state inference
% undefined.
%
% \paragraph{Boundary preservation.}
% No method merges a primary input or primary output into a super-node with a different role.
% The PI/PO boundary is the circuit's fixed interface, and dissolving it changes what the graph
% represents rather than merely coarsening it.
%
% % TODO: super-node content is not yet exploited; do not write as though it is.
%
% \subsubsection{Domain-Specific Coarsening}
% \label{sec:method:summary:domain}
% The AIG-native slot. Two candidates are implemented; which one is run is not yet decided,
% and the decision waits on measured compression and retention on the real corpus.
%
% \begin{itemize}
%     \item \textbf{Cone coarsening.} Merges along two independent axes: chains of
%     single-fanout gates contract into one super-node, up to a capped chain length, reducing
%     circuit \emph{depth}; and gates within a bounded band of topological levels that share
%     a common immediate \emph{post}-dominator merge together, compressing parallel width
%     while leaving critical-path depth untouched, so the level-based positional encoding
%     stays exact. The band width and the chain cap are the only compression controls any
%     domain-specific candidate offers.
%     \item \textbf{MFFC skeleton coarsening.} Contracts each \emph{maximal fanout-free cone}
%     into one super-node: every AND gate whose fanout lies entirely inside one cluster joins
%     that cluster, iterated to a fixed point. What survives is the multi-fanout sharing
%     skeleton plus the fixed PI/PO interface. It is parameter-free, so its compression is a
%     property of the circuit and cannot be dialled.
% \end{itemize}
%
% The two are related rather than rival: chain contraction is the capped, chain-shaped special
% case of MFFC absorption, which additionally captures cones whose interior reconverges.
%
% MFFC coarsening invites an obvious objection, and the answer is the thesis claim in
% miniature. MFFCs are prime targets for synthesis rewriting, since refactoring collapses
% and re-expresses one MFFC at a time, so merging them might be expected to destroy exactly
% the signal the model must read. The claim here is the opposite, and it is falsifiable:
% intra-cone wiring is what local rewriting is \emph{free to change}, so it is noise with
% respect to the label, whereas the sharing structure that \emph{constrains} what rewriting can
% achieve is what survives the merge.
%
% Both candidates preserve acyclicity by the same argument: a cluster is absorbed only into the
% single cluster receiving all of its outgoing edges, which cannot close a cycle in a DAG\@. For
% cone coarsening this holds at band width zero, where the width axis merges only within a
% level; wider bands give up both the guarantee and the exactness of the level encoding.
%
% % TODO: report that the width axis must use post-dominators, if this method is written up.
% % TODO: neither cone nor MFFC compression is measured on the corpus.
%
% \subsubsection{Graded Colour Refinement}
% \label{sec:method:summary:wl}
% The exact anchor of the study. Nodes are merged when they agree under $d$ rounds of colour
% refinement, with $d$ set equal to the encoder's message-passing depth so that merged nodes
% are exactly those the network cannot distinguish. A count-cap $c$ interpolates between two
% named endpoints \citep{bollen2023exact}: at $c = \infty$ refinement distinguishes nodes by
% the full multiset of neighbour colours and the merge is \emph{lossless} for a summing GNN\@;
% at $c = 1$ it sees only the set of neighbour colours, which is bisimulation: coarser,
% cheaper, and lossy for a counting model.
%
% Two AIG adaptations are required. Node colours are initialised from the four-way type and
% the level encoding rather than uniformly, so nodes at different depths never share a class
% and the pooled positional encoding of a super-node is exact. \textbf{Edge inversion is
% treated as a distinct relation} so that a normal and an inverted connection never
% contribute to the same colour class. Without this, logically distinct subgraphs merge. Refinement runs backward from the primary outputs
% by default; direction is a parameter, and it changes both how much compresses and which structure
% survives.
%
% At $c = \infty$ this method is orthogonal to the optimization being predicted by
% construction: it removes only what the model provably cannot see, so it cannot erase the
% label. That is what makes it the positive control for RQ5 rather than merely another
% contender, and realising the guarantee end to end requires the schema
% of~\ref{sec:method:architecture:exact}.
%
% The obvious risk is that structural hashing has already removed the easy equivalences
% (\ref{sec:prelim:algorithms:primitives}), since every corpus graph is strashed. What this
% method finds is therefore \emph{residual} redundancy beyond one hop, and how much of it
% exists is itself a reportable statistic about AIGs.
%
% % TODO: run the residual-redundancy probe over depths 1 to 4; depth 1 subsumes k-SNAP \citep{tian2008ksnap}.
%
% \subsubsection{Convolution Matching}
% \label{sec:method:summary:convmatch}
% The strongest general-purpose competitor: ConvMatch \citep{dickens2024convmatch} merges nodes
% that are equivalent with respect to the \emph{graph convolution operation} itself rather than
% with respect to a graph property. As for every other method, merges never cross node types.
% The restriction is imposed on the candidate set, so the convolution objective itself is
% unmodified.
%
% It is the bar the domain-specific method must clear, and it is the right bar precisely
% because it is GNN-aware but domain-blind: if an AIG-specific method beats it, the domain
% claim is earned against a method optimising the same objective, not against a strawman.
%
% Its behaviour on a deep DAG with inverted edges is untested in the literature, and one
% property found here explains why it may differ from the published setting. Its merge cost is
% an unweighted sum over the convolution output, so merging two low-degree nodes scores better
% than merging two genuine convolution twins of higher degree, regardless of similarity.
% Convolution equivalence therefore decides \emph{between} pairs of comparable degree, not
% across them, which on a graph with the degree profile of an AIG biases merging toward the
% input frontier.
%
% % TODO: measure the per-graph cost of probe sampling before quoting it.
%
% \subsubsection{Random Within-Type Merging}
% \label{sec:method:summary:random}
% Uniformly random nodes of the same type are merged until the node count reaches a target.
% Merges are restricted to one node type so that the comparison isolates \emph{which} nodes are
% merged rather than confounding it with whether the interface survives, matching the boundary
% guarantee of every other method.
%
% This serves two purposes, and the second is the important one.
%
% As a \textbf{naive floor}, it answers the question a reader will ask first: if a
% sophisticated method does not beat random merging at the same compression, its sophistication
% bought nothing. It is also the exact counterpart of random edge dropout in the sparsification
% family, so both halves of the study have a floor constructed the same way.
%
% As a \textbf{matching instrument}, it is what makes RQ4 answerable at all. Compression is a
% free parameter for only some methods: ConvMatch takes a target ratio directly, cone
% coarsening offers only coarse integer knobs, and MFFC coarsening and colour refinement have
% no compression parameter whatsoever. There is therefore no general way to make two \emph{real}
% methods meet at the same ratio. Random merging reaches any target exactly by construction, so
% each method is instead compared against a random control run at that method's own achieved
% compression. The paired design replaces a matching problem that has no solution.
%
% % TODO: not built. The cost is the training runs, one per matched point, not the code.
% \todo{completely look at and redo by hand comment out all this text and lets start with the basic concept i am plannign to implement. keep the text but commented out}
